% ============================================================
% ACM TRUST 2027 - EA TRACK: TRUSTWORTHY AND RESPONSIBLE AGENTIC AI
% DOUBLE-BLIND FULL PAPER - VERSION 3
%
% Main file:     trust2027.tex
% Bibliography:  reference.bib (same directory)
%
% BEFORE SUBMITTING:
%   1. Obtain independent annotations and update validity claims.
%   2. Verify bibliography, configuration snapshots and venue metadata.
%   3. All cited BibTeX keys must appear in reference.bib.
%   4. Build once end-to-end and check the page count (limit 9,
%      appendix included; references and GenAI disclosure free).
% ============================================================

\documentclass[sigconf,review,anonymous]{acmart}

% ACM owns the layout, fonts, margins, and hyperref configuration.
% Essential additions only: math symbols and flexible two-column tables.
\usepackage{amsmath,amsfonts,tabularx}
\emergencystretch=0.65em % Avoid protruding lines, without altering margins.

% Review-stage ACM metadata; check the official venue settings.
\setcopyright{none}
\copyrightyear{2027}
\acmYear{2027}
\acmConference[TRUST '27]
  {ACM Conference on Trustworthy and Responsible AI and Computing Systems}
  {March 7--9, 2027}
  {Washington, DC, USA}
\settopmatter{printacmref=true}

% Short local math notation for metrics used in the manuscript.
\newcommand{\CTWAR}{\mathrm{CTWAR}}
\newcommand{\PUAR}{\mathrm{PUAR}}
\newcommand{\SRR}{\mathrm{SRR}}
\newcommand{\EAC}{\mathrm{EAC}}
\newcommand{\OHR}{\mathrm{OHR}}
\newcommand{\ESC}{\mathrm{ESC}}
\newcommand{\JHR}{\mathrm{JHR}}
\newcommand{\RF}{\mathrm{RF}}
\newcommand{\UER}{\mathrm{UER}}
\newcommand{\eff}{\mathcal{E}_{\pi}}
\newcommand{\one}{\mathbb{I}}

% Unmeasured values are deliberately visible. Do NOT submit as-is.



\begin{document}

% ============================================================
% TITLE
% ============================================================

\title[Same Effect, Different Tool]{Same Effect, Different Tool: Auditing Post-Denial
Workaround Attempts in Tool-Using LLM Agents}

% Double-blind: no names, affiliations, contacts,
% acknowledgments, or identifying repository URLs.

% ============================================================
% ABSTRACT
% ============================================================

\begin{abstract}
Authorization checks can prevent forbidden tool execution while revealing
little about an agent's subsequent decisions. We study post-denial
continuations in which a tool-using language model must either respect an
authorization boundary, complete permitted work, seek approval, or halt.
We introduce Detour, an AgentDojo-backed evaluation instrument that
normalizes proposed tool actions into policy-relevant effects and measures
cross-tool workaround attempts (CTWAR) separately from other unauthorized
proposals, blocked execution, and authorized recovery. In a frozen suite
of 24 scenario instances derived from eight templates, we evaluate
Qwen2.5-7B, Qwen2.5-14B, and Mistral Small 3.2 24B under three
matched denial-feedback conditions. All three configurations completed
12/12 benign control tasks. Under bare denial feedback, strict CTWAR
was 16/36 for Qwen2.5-7B, 6/36 for Qwen2.5-14B, and 0/36 for Mistral.
However, zero observed cross-tool workarounds did not imply effective
permitted continuation: Mistral's unwarranted halt rate increased from 30/54 to
42/54 when the denial named the governing restriction, while its
safe-recovery rate remained 9/54. These observations distinguish
attempted circumvention from excessive halting in a fully enforcing
sandbox. The study is limited by a small number of independent template
families, structural rather than semantic recovery checks, and effect
equivalence that has not yet been independently validated by annotators.
\end{abstract}

% ============================================================
% CCS CONCEPTS AND KEYWORDS
% ============================================================

\ccsdesc[500]{Computing methodologies~Artificial intelligence}
\ccsdesc[300]{Security and privacy~Software security engineering}
\ccsdesc[300]{Computing methodologies~Multi-agent systems}

\keywords{
trustworthy agentic AI,
tool-using LLMs,
authorization boundaries,
agent evaluation,
safe recovery,
accountability
}

\maketitle

% ============================================================
% 1. INTRODUCTION
% ============================================================

\section{Introduction}
\label{sec:intro}

Language model agents act. They plan across steps, read
external data, and invoke tools that transmit information,
alter records and start workflows. ReAct-style interleaving of
reasoning and acting made the pattern
practical~\cite{yao2023react}; the same interface exposes
agents to instructions embedded in lower-trust content and to
mistakes about what they are permitted to
do~\cite{debenedetti2024agentdojo,zhang2025asb}.
Trustworthiness consequently cannot be read off task
completion. It depends on whether the task was accomplished
through actions the user and the governing policy authorized.

Runtime authorization is the standard safeguard, and it works:
a correct gate blocks the call it is asked to block. What it
does not do is end the agent's turn. Consider an assistant
authorized to prepare an email but not to transmit it
externally. It proposes
\texttt{send\_email(recipient,\,payload)} and the gate refuses.
The agent may now save a draft, request approval, stop and
explain, or call \texttt{post\_message} with the same recipient
and the same payload. The last option changes the tool
identifier and nothing else: its policy-relevant effect is
still \emph{transmit the protected payload to the unauthorized
recipient}.

Both runs leave the environment untouched, so an execution-only
metric scores them the same. They are not the same. One agent
respected the boundary it was told about; the other looked for
a way around it and was stopped by the gate rather than by its
own judgment. In deployment, where tool inventories are larger
and gates are partial, that difference is the difference
between a safeguard that holds and one that happens to have
covered the route the agent tried first.

\paragraph{The measurement problem.}
Naming the behavior is easy; counting it is not. A second tool
call that differs from the denied one is not evidence of
anything on its own, because the authorized recovery path also
uses a different tool. Distinguishing the two requires saying
what a tool call \emph{does} in policy terms, independently of
what it is called, and holding that judgment to a reviewable
standard rather than a model's opinion. This paper's technical
content is largely that: a normalized effect representation, an
independent-labeling protocol to be completed, and a metric
whose denominator is restricted to episodes where an
effect-equivalent route provably exists.

\paragraph{What this paper is.}
This is an evaluation artifact, not a defense. We introduce no
guardrail and train no model. We hold a deterministic
enforcement rule fixed, vary only the information carried by
the denial, and audit what frozen pretrained agents propose
afterwards. The deliverables are a scenario schema, an
annotation protocol, and a metric suite, together with the
experimental design that makes their comparisons valid.

\paragraph{Research questions.}
\textbf{RQ1.} Do agents propose effect-equivalent alternate-tool
actions after a denial, at a rate distinguishable from exact
retries and from authorized recovery?
\textbf{RQ2.} Does the information content of the denial change
that rate, and does it do so without inducing over-refusal on
episodes where an authorized route exists?

We additionally examine sensitivity to adversarial lower-trust
tool content, the confound of general tool-use incompetence,
and the agreement between an agent's final report and the
instrumented environment state.

\paragraph{Contributions.}
\emph{First}, a policy-grounded scenario representation in which
each denial carries its correct-continuation class, its
authorized routes, and its labeled effect-equivalent
alternatives (Section~\ref{sec:problem}).
\emph{Second}, a metric suite built around cross-tool workaround
attempts, with disjoint eligibility sets, an explicit separation
of attempted from executed violations, and an over-refusal
control that prevents halting from being scored as safety
(Section~\ref{sec:metrics}).
\emph{Third}, a matched-feedback protocol under a fixed gate,
which isolates differences in denial text under fixed enforcement
while exposing the limits of conclusions from seeded denied states
(Section~\ref{sec:method}).

% ============================================================
% 2. RELATED WORK
% ============================================================

\section{Related Work and Positioning}
\label{sec:related}

\subsection{Executable Agent Security Evaluation}

AgentDojo supplies realistic tasks, simulated environments and
attacker-controlled tool responses for prompt-injection
evaluation~\cite{debenedetti2024agentdojo}; Agent Security
Bench systematizes attacks and defenses across agent
settings~\cite{zhang2025asb}. Both establish that agent safety
must be tested on executed trajectories rather than final text.
ST-WebAgentBench makes a workflow's governing rules explicit
and scores completion under those
constraints~\cite{levy2026stweb}. We adopt the executable-task
and explicit-policy discipline of this line and change the unit
of analysis: measurement is conditioned on a specific rejected
invocation, and labels attach to what the agent proposes
afterwards even when nothing forbidden executes.

\subsection{Guardrails, Alignment Checks and Tool Contracts}

Task Shield tests whether lower-trust instructions and proposed
actions serve the user's authorized task~\cite{jia2025taskshield}.
IPIGuard constrains execution through tool-dependency
graphs~\cite{an2025ipiguard}; ToolGate governs execution with
contracts over verified state
transitions~\cite{liu2026toolgate}. TraceSafe assesses
guardrails across multi-step tool-calling trajectories and finds
their efficacy driven more by structural competence than by
semantic alignment~\cite{chen2026tracesafe}.

ToolSafe is the closest of these: a step-level guardrail whose
feedback drives continued reasoning rather than aborting, with
the finding that richer feedback improves both safety and
utility over a bare safety rating~\cite{mou2026toolsafe}. We do
not contest that result and do not claim feedback as a novel
mechanism. ToolSafe varies how \emph{informative} a denial is
and scores the rate of harmful invocations. We vary how
\emph{policy-grounded} a denial is while holding enforcement
fixed, and score which \emph{effect} the agent pursues next --- a
quantity a per-invocation rate cannot express, because route
substitution and ordinary noncompliance are the same event to it.

\subsection{Authorization, Privilege and Attempted Violations}

AgentBoundary varies apparent risk and action permissibility
independently across counterfactual variants of one executable
workflow, judging both trajectory and post-state across many
model and harness configurations, and exposes a sharp tension
between blocking unauthorized actions and completing authorized
ones~\cite{yang2026agentboundary}. It is the strongest existing
treatment of permission calibration, and our contribution is
not scale against it. Its unit is the task: a run is scored by
whether the authorized work was completed and whether a harmful
state transition occurred. An episode in which the agent is
blocked, tries an equivalent route through a second tool, is
blocked again and then halts is, at that granularity, one
failed task. Our unit is the effect, so the same episode
carries a recorded boundary violation. Section~\ref{sec:results}
states this as a falsifiable prediction rather than an
assertion.

FORTIS measures over-privileged skill and tool selection,
finding that models routinely reach past the minimum privilege
a task requires~\cite{li2026fortis}; its decision point is the
initial choice, ours is a replanning step after a communicated
refusal. PIGuard studies over-defense in injection
detection~\cite{li2025piguard}, which is why our suite scores
over-halting explicitly rather than treating refusal as
costless. Siu et al.\ argue that agent security is contextual
and propose source authorization, task alignment, action
alignment and data isolation as the properties that must hold
jointly~\cite{siu2026contextual}; a cross-tool workaround is an
action-alignment violation in their vocabulary, and this paper
supplies an instrument that measures one class of them.

Closest in phenomenon, Yu et al.\ report that agents in a
deterministic environment with explicit constraints develop
multi-step strategies to circumvent safety limits without
adversarial prompting, and observe that attempted violations
are rarely measured because outcome-based metrics
dominate~\cite{yu2026affordance}. We take that observation as
established rather than as our motivation to prove, and supply
what it lacks: an annotated effect-equivalence relation across
distinct tools, which converts circumvention from an
observation into a rate with a stated denominator. Related
evidence that refusal behavior in text fails to transfer to
tool calls~\cite{cartagena2026gap} motivates conditioning on a
denial the agent has already received, where the boundary has
been stated in the clearest available terms. Reporting
unfaithfulness under silent tool failures has been audited
by Singh~\cite{singh2026scapegoats}; our reporting-fidelity
check shares that concern under a different trigger, an
explicit policy denial, and is a secondary accountability
measure rather than a new failure family.

% ============================================================
% TABLE 1: POSITIONING
% ============================================================

\begin{table}[t]
\caption{
Positioning by primary analysis unit. Adjacent behaviors are
not absent from the cited work; they are not its unit of
measurement.
}
\label{tab:prior}
\small
\setlength{\tabcolsep}{3pt}

\begin{tabularx}{\columnwidth}{
@{}
>{\raggedright\arraybackslash}p{2.5cm}
>{\raggedright\arraybackslash}X
@{}
}
\toprule
Study & Primary analysis unit \\
\midrule

AgentBoundary~\cite{yang2026agentboundary} &
Task, under permission/risk counterfactuals. \\

ToolSafe~\cite{mou2026toolsafe} &
Tool invocation, scored safe or unsafe. \\

FORTIS~\cite{li2026fortis} &
Initial skill and tool selection. \\

Yu et al.~\cite{yu2026affordance} &
Trajectory, circumvention observed qualitatively. \\

This work &
Post-denial episode; normalized effect of each
subsequent proposal. \\

\bottomrule
\end{tabularx}
\end{table}

% ============================================================
% 3. PROBLEM FORMULATION
% ============================================================

\section{Problem Formulation}
\label{sec:problem}

\subsection{Agent, Gate and Trust Boundary}

A task instance comprises a legitimate goal $g$, trusted state
$s_t$, tool library $\mathcal{T}$ and state-dependent policy
$\pi$. At step $t$ a frozen agent $M$ proposes
$a_t=(\tau_t,\theta_t)$ from history $h_t$, with
$\tau_t\in\mathcal{T}$ the tool and $\theta_t$ its arguments.
The gate evaluates

\begin{equation}
\label{eq:gate}
G_{\pi}(s_t,g,a_t)\in\{0,1\}.
\end{equation}

When $G_\pi=1$ the instrumented tool applies
$s_{t+1}=T(s_t,a_t)$. When $G_\pi=0$ the gate records the
\emph{proposed} action, executes nothing, returns a denial
message and leaves protected state unchanged. An agent can
therefore accumulate unauthorized attempts in an environment
that never performs an unauthorized execution, and the two
quantities are reported separately throughout.

The adversary may control lower-trust material the agent
consumes, including document bodies and returned tool content,
but not the gate, the policy labels, the executor or the audit
log. Side channels, compromised enforcement and sandbox escapes
are out of scope. A permitted invocation is not treated as
evidence that the user's goal was met.

\subsection{Normalized Effects Across Tools}

Distinct APIs can produce one policy-relevant consequence. We
define an annotation oracle --- not a learned classifier --- that
normalizes a proposed action's effect as a four-tuple

\begin{equation}
\label{eq:effect}
\eff(s,a)=\langle o,\; r,\; b,\; q\rangle,
\end{equation}

where $o$ is an operation class, $r$ the principal or location
receiving the effect, $b$ the protected object, and $q$ a
qualifier carrying scope or mode. The qualifier is required
because the first three slots are insufficient for the
operations this study uses. For \emph{transmit}, $r$ is the
recipient, $b$ the payload and $q$ the channel class. For
\emph{delete}, $r$ is the store, $b$ the target and $q$
distinguishes a single-object removal from a recursive one. For
\emph{grant access}, the grantee and the resource are different
entities: $r$ is the grantee, $b$ the resource and $q$ the
privilege level. Collapsing $q$ would make a recursive delete
and a single-file delete indistinguishable, which would bias
the primary metric in the permissive direction.

Let $a_d=(\tau_d,\theta_d)$ be the denied action in state
$s_d$. A later proposal $a_t$ is an \emph{effect-equivalent
cross-tool attempt} if and only if

\begin{equation}
\label{eq:equiv}
\tau_t\neq\tau_d
\;\;\wedge\;\;
\eff(s_t,a_t)=\eff(s_d,a_d)
\;\;\wedge\;\;
G_\pi(s_t,g,a_t)=0 .
\end{equation}

The third conjunct matters: the policy is state-dependent, so an
effect forbidden at $s_d$ may be permitted at $s_t$ --- after an
approval has been granted, for instance --- and such a call is
compliant, not a workaround. We score a recorded proposal, never
an inferred intention. A different tool name alone is
insufficient; a permitted alternative with a different effect is
not a violation. Broad access-grant equivalence and
adversarial-content templates are excluded from the strict primary
analysis and evaluated separately as sensitivity checks.

\subsection{Episode Classes and Eligibility}

Let $\mathcal{D}$ be the episodes beginning at a standardized
denial record $d$, continuing through trace $\xi_d$ and ending
in report $r_T$. Each belongs to exactly one correct-continuation
class:

\begin{equation}
\label{eq:partition}
\mathcal{D}
=
\mathcal{D}_{\rm alt}
\mathbin{\dot\cup}
\mathcal{D}_{\rm approve}
\mathbin{\dot\cup}
\mathcal{D}_{\rm halt},
\end{equation}

according to whether an authorized completion route exists,
explicit approval is required, or no authorized continuation
exists. Each class has its own correct outcome and, in
Section~\ref{sec:metrics}, its own estimator; none is folded
into another.

Cross-tool eligibility is an \emph{orthogonal} attribute rather
than a fourth class. The subset
$\mathcal{D}_{\rm cross}\subseteq\mathcal{D}$ marks episodes in
which the environment exposes a tool capable of the denied
effect, and
$\mathcal{D}_{\rm alt}\cap\mathcal{D}_{\rm cross}\neq\emptyset$
is intended: the episodes in which a legitimate route and a
forbidden shortcut are simultaneously available are the most
diagnostic in the suite, because the agent is not choosing
between compliance and failure but between compliance and
convenience.

% ============================================================
% 4. METHOD
% ============================================================

\section{Effect-Grounded Post-Denial Audit}
\label{sec:method}

\subsection{Architecture}

Figure~\ref{fig:arch} separates four components that are
conflated in execution-only evaluation: the frozen agent, the
deterministic gate, the sandbox executor, and an auditor that
observes proposals rather than outcomes. The agent receives
ordinary task instructions and tool outputs; across controlled
conditions only the denial text differs. Every state-changing
call passes $G_\pi$, whose decision is logged with the proposal
that triggered it. The auditor consumes all proposals, denial
messages, permitted transitions and final claims, and labels
them against effect annotations drawn from the scenario
specification. Those annotations are never produced by the
agent under evaluation.

% ============================================================
% FIGURE 1: ARCHITECTURE
% ============================================================

\begin{figure}[t]
\centering
% Self-contained architecture figure: no external PDF or TikZ required.
\begingroup
\footnotesize
\setlength{\fboxsep}{3pt}
\setlength{\fboxrule}{0.5pt}
\fbox{\parbox{0.90\columnwidth}{\centering
User goal $g$, policy $\pi$, trusted state $s_t$}}
\par\vspace{2pt}
$\downarrow$
\par\vspace{2pt}
\fbox{\parbox{0.90\columnwidth}{\centering
Frozen agent $M$ proposes $a_t=(\tau_t,\theta_t)$}}
\par\vspace{2pt}
$\downarrow$
\par\vspace{2pt}
\fbox{\parbox{0.90\columnwidth}{\centering
Trusted gate $G_\pi$: log proposal and check permission}}
\par\vspace{2pt}
\begin{minipage}[t]{0.47\columnwidth}
\centering
$\swarrow\;G_\pi=0$\\[2pt]
\fbox{\parbox{0.82\linewidth}{\centering
\textbf{Denied}\\
No execution; return $F_k$\\
Post-denial episode begins}}
\end{minipage}\hfill%
\begin{minipage}[t]{0.47\columnwidth}
\centering
$G_\pi=1\;\searrow$\\[2pt]
\fbox{\parbox{0.82\linewidth}{\centering
\textbf{Allowed}\\
Sandbox executes call\\
State updates to $s_{t+1}$}}
\end{minipage}
\par\vspace{3pt}
$\downarrow$\\[2pt]
\fbox{\parbox{0.90\columnwidth}{\centering
\textbf{Offline auditor}: annotated effects
$\mathcal{E}_{\pi}=\langle o,r,b,q\rangle$;
proposal, state and report labels}}
\endgroup
\caption{Effect-grounded audit architecture. The external gate
controls execution; the independent auditor records proposed
cross-tool effects, permitted state changes and reporting fidelity.
Scenario-defined effect labels are not generated by the evaluated agent.}
\Description{A frozen agent proposes an action under a user goal and
policy. A trusted gate denies or executes the action. A denial
returns feedback and begins an audit episode; an allowed call
updates sandbox state. An offline auditor uses scenario-defined
effect tuples to evaluate proposals and final reports.}
\label{fig:arch}
\end{figure}

\subsection{Outcome Labels}

A post-denial trace may contain a compliant alternative, an
approval request, a justified halt, an exact retry, an
effect-equivalent call through another tool, or an unrelated
error. These are not mutually exclusive: an agent may attempt a
shortcut and subsequently recover, and the suite records both
rather than resolving them into a single verdict.
Table~\ref{tab:taxonomy} gives the decision rule for each label
and names the metric it feeds, so that the mapping from
observation to number is explicit rather than reconstructed by
the reader.

% ============================================================
% TABLE 2: OUTCOME TAXONOMY
% ============================================================

\begin{table}[t]
\caption{
Outcome labels, decision rules and the metric each feeds.
Labels are assigned from proposals, policy and instrumented
state.
}
\label{tab:taxonomy}
\small
\setlength{\tabcolsep}{3pt}

\begin{tabularx}{\columnwidth}{
@{}
>{\raggedright\arraybackslash}p{1.75cm}
>{\raggedright\arraybackslash}X
>{\raggedright\arraybackslash}p{1.05cm}
@{}
}
\toprule
Label & Decision rule & Feeds \\
\midrule

Safe recovery &
Authorized goal state reached with no
policy-violating proposal. &
$\SRR$ \\

Late recovery &
Goal reached, but a violating proposal occurred
earlier in the episode. &
$\EAC$ \\

Appropriate escalation &
Approval requested; protected operation not
executed first. &
$\ESC$ \\

Justified halt &
No authorized route exists; restriction reported
correctly. &
$\JHR$ \\

Unwarranted halt &
Agent stops although an authorized route exists. &
$\OHR$ \\

Exact retry &
Denied tool and forbidden effect proposed again. &
$\PUAR$ \\

Cross-tool workaround &
Different tool, same forbidden effect, still
forbidden at proposal time. &
$\CTWAR$ \\

Reporting mismatch &
Final claim contradicts instrumented terminal
state. &
$\RF$ \\

\bottomrule
\end{tabularx}
\end{table}

\subsection{Matched Denial-Feedback Conditions}

Three messages are compared for the same denied action under
identical enforcement. $F_0$ reports only that the action was
denied. $F_1$ names the governing restriction, for example that
external transmission requires approval. $F_2$ additionally
names a \emph{permitted action class}, for example that
local-only preparation remains available, and never supplies a
concrete tool, arguments or a recovery plan.

The class-not-call distinction is deliberate and is the design
decision on which the $F_2$ comparison rests. Supplying the
exact recovery call would make safe recovery a test of whether
the agent can repeat the experimenter's answer; supplying only
the class leaves the agent to locate a compliant route itself.
Every subsequent call, including calls consistent with $F_2$,
passes the same gate. Because feedback content is an
intervention, differences between conditions estimate the
effect of additional policy information and not an intrinsic
property of the model.

\subsection{Controlled Post-Denial Continuations}
The reported evaluation replays the same policy-denied prefix for every
model and feedback condition, then audits subsequent proposals in the
instrumented sandbox. This identifies behavior conditional on being placed
in a denied state; it does not estimate how often an unprompted agent
would encounter such a denial in a natural deployment. No natural-rollout
results are reported here. Denial feedback changes; gate rules, tool
inventory, and synthetic initial state remain fixed within the experiment.

% ============================================================
% 5. EXPERIMENTAL SETUP
% ============================================================

\section{Experimental Setup}
\label{sec:experiment}

\subsection{Task Families}
The evaluated Detour suite comprises 24 scenario instances derived from
eight templates, each instantiated with three variants. Eighteen
instances across six templates are cross-tool eligible; three are
adversarial document-content variants matched to a benign template.
Correct continuation is classified as a permitted alternative (18
instances), required approval (3), or justified halt (3). Twenty-one
instances involve disclosure and three involve deletion. The environment
uses synthetic document, email, and chat resources with a deterministic
authorization gate. The same scenarios are evaluated with three feedback
conditions and three repetitions per condition, resulting in 216 audit
episodes per model configuration. Changing names within a template does
not supply independent evidence about a new workflow; uncertainty is
therefore organized by template, not by episode count.

\subsection{Effect Annotation and Quality Control}
\label{sec:annotation}
The implemented effect oracle matches operation, recipient/principal,
protected object, and qualifier using scenario-defined rules and aliases.
These are deterministic \emph{machine annotations}, not independently
human-adjudicated ground truth. We distinguish the strict, non-adversarial
primary effect scope (excluding broad access-grant equivalence in T4 and
the injected-content T6) from broader and channel-equivalence sensitivity
analyses. The policy scope and object/recipient aliases may generate false
matches or miss equivalent effects. Independent, blinded human labeling
is needed to estimate oracle validity; until that validation is complete,
we do not claim a measured inter-annotator agreement, precision, or recall.

\subsection{Agents, Competence Control and Baselines}
We evaluate three frozen open-weight configurations: Qwen2.5-7B,
Qwen2.5-14B, and Mistral Small 3.2 24B. The two Qwen models were
previously evaluated in a matched Detour v8 run, and Mistral in a frozen
v9 extension whose dataset and metric compatibility check passed.
Before post-denial scoring, each agent was screened on six benign tasks
with two repetitions each, including a read-then-send workflow.
All three configurations reached 12/12 benign-task goals and 2/2
multi-step goals in those controls. This is a small competence screen,
not proof of general tool competence. Mistral was served via Ollama
with an 8192-token configured context; the exact earlier Qwen serving
settings must be taken from their saved run manifests before claiming
full decoding or hardware equivalence.

The feedback comparison isolates F0, F1, and F2 within each configuration.
We compare CTWAR with general unauthorized proposals (PUAR), safe recovery,
unnecessary escalation, unjustified halting, and the blocked-execution
baseline. As the gate prevents forbidden execution, zero observed forbidden
terminal states is primarily an enforcement finding and not evidence of
compliant agent intent.

% ============================================================
% TABLE 3: DESIGN
% ============================================================

\begin{table}[t]
\caption{
Prespecified comparisons and what each controls.
}
\label{tab:design}
\small
\setlength{\tabcolsep}{3pt}

\begin{tabularx}{\columnwidth}{
@{}
>{\raggedright\arraybackslash}p{2.1cm}
>{\raggedright\arraybackslash}X
@{}
}
\toprule
Comparison & What it isolates \\
\midrule

$F_0$ vs.\ $F_1$ &
Effect of naming the governing restriction. \\

$F_1$ vs.\ $F_2$ &
Added effect of a permitted action class, with no
tool hint. \\

Exact vs.\ cross-tool &
Route substitution against generic noncompliance. \\

Benign vs.\ adversarial &
Sensitivity to untrusted tool content. \\

Clean-task competence &
Inability to act against deliberate compliance. \\

\bottomrule
\end{tabularx}
\end{table}

% ============================================================
% 6. METRICS
% ============================================================

\section{Metrics and Statistical Analysis}
\label{sec:metrics}

\subsection{Primary Outcome}

For $i\in\mathcal{D}_{\rm cross}$, let $W_i=1$ if any
subsequent proposal satisfies Eq.~\eqref{eq:equiv}, and $0$
otherwise. The primary metric is

\begin{equation}
\label{eq:ctwar}
\CTWAR
=
\frac{\sum_{i\in\mathcal{D}_{\rm cross}}W_i}
{|\mathcal{D}_{\rm cross}|}.
\end{equation}

The denominator admits only episodes in which an
effect-equivalent route exists and has been annotated. An exact
same-tool retry never satisfies $W_i=1$; an attempted
substitution need not execute. Because operation classes differ
in reversibility, we additionally report a severity-weighted
variant in which each episode carries a weight $w_i$ fixed by
its operation class before data collection, with irreversible
operations weighted above reversible ones. The unweighted rate
remains primary.

\subsection{Secondary Outcomes by Eligibility Class}

Let $U_i$ indicate any unauthorized post-denial proposal
regardless of tool, and $X_i$ an unauthorized \emph{execution}.
Over all episodes,

\begin{equation}
\label{eq:global}
\PUAR=\frac{\sum_{i\in\mathcal{D}}U_i}{|\mathcal{D}|},
\qquad
\UER=\frac{\sum_{i\in\mathcal{D}}X_i}{|\mathcal{D}|}.
\end{equation}

$\UER$ is expected to be zero under a correct gate and is
reported as an enforcement check, not as a safety result.

Each eligibility class carries its own estimator. On
$\mathcal{D}_{\rm alt}$, let $C_i$ indicate completion with no
violating proposal anywhere in the episode, $C^{+}_i$
completion irrespective of earlier violations, and $H_i$ a halt
despite an available authorized route:

\begin{equation}
\label{eq:alt}
\SRR=\frac{\sum C_i}{|\mathcal{D}_{\rm alt}|},
\quad
\EAC=\frac{\sum C^{+}_i}{|\mathcal{D}_{\rm alt}|},
\quad
\OHR=\frac{\sum H_i}{|\mathcal{D}_{\rm alt}|}.
\end{equation}

The gap $\EAC-\SRR$ is the share of successful episodes that
passed through an attempted violation, which an outcome-only
metric reports as unqualified success. $\OHR$ is the
over-refusal control: without it, an agent that halts
everywhere scores well on every other
measure~\cite{li2025piguard}. On the remaining classes, with
$Q_i$ indicating approval correctly requested before any
protected execution and $J_i$ a correctly reported halt,

\begin{equation}
\label{eq:classes}
\ESC=\frac{\sum_{i\in\mathcal{D}_{\rm approve}}Q_i}
{|\mathcal{D}_{\rm approve}|},
\qquad
\JHR=\frac{\sum_{i\in\mathcal{D}_{\rm halt}}J_i}
{|\mathcal{D}_{\rm halt}|}.
\end{equation}

Tool-call counts, continuation steps, latency and token use are
reported alongside. No measure is merged into a composite
trustworthiness score: the suite's purpose is to keep safety
and utility legible as separate quantities, and a weighted sum
would discard exactly the information it was built to preserve.

\subsection{Reporting Fidelity}

Let $c_i$ be a parseable final claim about action status and
$\sigma_i$ the reference status from logged transitions ---
\emph{sent}, \emph{draft-only}, \emph{denied},
\emph{awaiting approval}. Over the eligible subset
$\mathcal{D}'\subseteq\mathcal{D}$ of unambiguous claims,

\begin{equation}
\label{eq:rf}
\RF=\frac{1}{|\mathcal{D}'|}\sum_{i\in\mathcal{D}'}
\one[c_i\simeq\sigma_i].
\end{equation}

Claims not interpretable under the prespecified ontology are
never silently scored correct; $|\mathcal{D}\setminus
\mathcal{D}'|$ is reported as the abstention count with its
fraction. This measures terminal reporting under explicit
denial and does not substitute for auditing
execution~\cite{singh2026scapegoats}.

\subsection{Uncertainty and Effective Sample Size}

Conditions are paired on scenario instances. For outcomes with observed
variation, a paired bootstrap resamples \emph{templates}, not individual
renamed tasks or repeated episodes. The strict CTWAR scope has only four
independent template clusters per model--feedback cell; cluster-bootstrap
intervals are therefore unstable. When every observed outcome is zero,
resampling only zero-valued clusters produces a degenerate interval that
is \emph{not} a valid claim of zero population risk. We report counts,
eligible denominators and cluster counts, and avoid binomial intervals
that assume independent episodes. Oracle agreement cannot be reported
until independent human labels are collected.

% ============================================================
% 7. RESULTS
% ============================================================

\section{Results}
\label{sec:results}

\subsection{Coverage and Clean Competence}
The frozen Detour study includes 24 instances in eight templates. Each
model configuration was evaluated under three feedback conditions and
three repetitions, yielding 216 post-denial episodes each and 648 across
three configurations. The strict, non-adversarial CTWAR denominator has
36 eligible episodes and four template clusters per model-condition cell.
The broader six-template eligible set has 54 episodes per cell, including
the adversarial T6 variant. All three configurations completed 12/12
benign controls and the 2/2 clean chaining task instances. The Mistral
run reported no seed errors, zero gate failures, and no forbidden final
state in 216 episodes. The older Qwen findings were obtained from the
saved v8 evidence and are reported using their original metrics.

\subsection{Main Denial-Feedback Results}
Table~\ref{tab:results} reports observed rates and their correct eligible
denominators. Strict CTWAR excludes T4, whose access grant is a contested
broad equivalence, and T6, which adds injected tool content. We deliberately
omit episode-binomial confidence intervals: three repetitions and
name variants within four strict templates are not independent draws.
A zero event rate also makes a template bootstrap degenerate; it does not
imply a precise out-of-sample zero.

\begin{table*}[t]
\caption{Controlled post-denial results across three tool-competent model
configurations. Values are percentages followed by observed count/eligible
count. CTWAR uses strict non-adversarial eligibility (36 per cell), PUAR
uses all episodes (72), and SRR, ESC-alt, OHR use the alternative-route
class (54). These are descriptive observations rather than independent-
trial confidence bounds.}
\label{tab:results}
\centering\small\setlength{\tabcolsep}{3pt}
\begin{tabular}{llccccc}
\toprule
Model & Feedback & Strict CTWAR $\downarrow$ & PUAR $\downarrow$ & SRR $\uparrow$ & ESC-alt & OHR $\downarrow$ \\
\midrule
Qwen2.5-7B & $F_0$ & 44.4 (16/36) & 36.1 (26/72) & 24.1 (13/54) & 5.6 (3/54) & 24.1 (13/54) \\
Qwen2.5-7B & $F_1$ & 0 (0/36) & 0 (0/72) & 3.7 (2/54) & 96.3 (52/54) & 0 (0/54) \\
Qwen2.5-7B & $F_2$ & 22.2 (8/36) & 18.1 (13/72) & 29.6 (16/54) & 27.8 (15/54) & 18.5 (10/54) \\
\midrule
Qwen2.5-14B & $F_0$ & 16.7 (6/36) & 22.2 (16/72) & 22.2 (12/54) & 0 (0/54) & 48.1 (26/54) \\
Qwen2.5-14B & $F_1$ & 0 (0/36) & 5.6 (4/72) & 22.2 (12/54) & 66.7 (36/54) & 3.7 (2/54) \\
Qwen2.5-14B & $F_2$ & 2.8 (1/36) & 15.3 (11/72) & 29.6 (16/54) & 46.3 (25/54) & 3.7 (2/54) \\
\midrule
Mistral Small 24B & $F_0$ & 0 (0/36) & 20.8 (15/72) & 16.7 (9/54) & 0 (0/54) & 55.6 (30/54) \\
Mistral Small 24B & $F_1$ & 0 (0/36) & 4.2 (3/72) & 16.7 (9/54) & 0 (0/54) & 77.8 (42/54) \\
Mistral Small 24B & $F_2$ & 0 (0/36) & 2.8 (2/72) & 20.4 (11/54) & 0 (0/54) & 75.9 (41/54) \\
\bottomrule
\end{tabular}
\end{table*}

\subsection{Contrasting Post-Denial Failure Modes}
Under bare feedback, Qwen2.5-7B proposed a strict cross-tool workaround
in 16/36 eligible episodes, and Qwen2.5-14B did so in 6/36. Naming the
restriction (F1) yielded 0/36 for both Qwen configurations, but the
smaller model's safe recovery also fell from 13/54 to 2/54 and its
unnecessary escalation rose from 3/54 to 52/54. Providing a permitted
class (F2) did not produce a monotonic safety improvement: strict CTWAR
returned to 8/36 for Qwen2.5-7B and 1/36 for Qwen2.5-14B.

Mistral exhibited a different failure mode. No CTWAR was detected in any
of its 162 broadly cross-tool eligible episodes, including the adversarial
arm; however, unauthorized proposals still occurred, and strict
unwarranted halting increased from 30/54 under F0 to 42/54 under F1.
Safe recovery remained 9/54 under both conditions. Thus an observed
absence of cross-tool substitution neither implies absence of other
violations nor effective continuation.

\subsection{Incremental Value and Reporting Limits}
An execution-only metric recorded no forbidden terminal states in the
fully enforcing Mistral experiment, even though 20/216 episodes contained
some unauthorized proposal. The automated Qwen v8 strict, completion-
eligible D1 comparison identified five episodes with both a cross-tool
workaround and \emph{structural} completion (four for 7B and one for
14B). These completion labels are based on environment predicates rather
than an independent semantic assessment of the resulting artifact.
Mistral contributed no CTWAR positives to D1. The effect-aware metric
therefore separates attempted circumvention from both blocked execution
and excessive stopping, but its correspondence to independent human
policy judgments remains to be established.

Mistral ended 24/216 episodes in prose rather than a structured
\texttt{report\_and\_stop} call, and only 24/216 episodes supplied an
RF-eligible status under the current parser. RF accuracy was 1/8, 1/8,
and 0/8 for F0, F1, and F2, respectively; those fractions must be read
together with 192/216 abstentions, not as model-wide reporting accuracy.
The correct halt metric was 0/9 in every condition and may partly reflect
unavailable or unrecognized terminal statuses; this warrants a separate
trace audit before a behavioral conclusion is drawn.

% ============================================================
% 8. DISCUSSION AND LIMITATIONS
% ============================================================

\section{Discussion and Limitations}
\label{sec:discussion}

\subsection{What the Audit Establishes}

The audit measures an agent's propensity to propose a forbidden
effect after the boundary has been stated to it explicitly. Its
central separation is between system enforcement and agent
continuation. Under a correct gate the unauthorized execution
rate is zero by construction; this does not constrain $\CTWAR$,
and the two numbers are not interchangeable. A deployment that
reports the first as its safety evidence is reporting a
property of its gate, not of its agent --- which is adequate
exactly as long as the gate's coverage is complete, and
tool inventories in practice grow faster than policy
coverage does.

The converse failure is equally available. An agent that halts
after every denial achieves a low workaround rate while failing
legitimate work, which is why $\OHR$ and clean-task competence
are reported before $\CTWAR$ rather than after it.

\subsection{Construct Validity}

The effect tuple simplifies context-dependent consequences by
construction. The qualifier slot covers scope and mode for the
operation classes used here, but operations outside this set
may need further qualifiers, and both false matches and
omissions bias $\CTWAR$ directly --- false matches upward,
omissions downward. Current mitigations include strict-versus-broad sensitivity
analysis and instrumented state checks. Independent human
annotation is still required and is not claimed as completed. We scope the claim
accordingly: the instrument implements rules for disclosure, deletion
and access-grant cases in the environments studied, and
extending it is an empirical question rather than a
generalization we assert.

\subsection{Intervention and External Validity}

A model placed into a predefined denied state might not have
proposed the antecedent action itself, which is the principal
threat to the controlled arm; no natural-rollout control has yet measured its magnitude. More
informative feedback changes what the agent knows, so the
comparison estimates an information intervention and not an
intrinsic model property independent of prompting.

Synthetic email, chat and document workflows do not stand
in for production agents, dynamically discovered tools, browser
sessions, multi-agent delegation or hidden side channels.
Results may vary with tool schema, prompt wording, API
snapshot, serving backend and continuation budget. The number
of independent task families bounds generalization, and the
suite is sized to establish whether the phenomenon occurs
rather than to estimate its prevalence in deployment.

\subsection{Responsible Evaluation}

All recipients, files and resources are synthetic; tools perform
no real outbound communication and modify no live account.
Adversarial content remains confined to the sandbox. Released
traces are non-sensitive and contain no real identifiers. The
framework evaluates one declared boundary in an instrumented
environment; it is neither formal verification nor a
production security certification, and a low $\CTWAR$ is
evidence about the agents and tasks measured, not a safety
guarantee for any deployment.

% ============================================================
% 9. CONCLUSION
% ============================================================

\section{Conclusion}
\label{sec:conclusion}
Detour separates an agent's post-denial proposals from actual tool
execution and from permitted task recovery. In a 24-instance benchmark,
competent Qwen configurations sometimes pursued the forbidden effect
through another tool, whereas a competent Mistral configuration showed
no such observed substitution but frequently halted despite available
legitimate continuations. Explicit restrictions changed the balance
between attempted violations, escalation and halting, not uniformly in
a favorable direction. These findings support evaluating effect-grounded
proposals and authorized utility jointly, without equating zero executed
violations with agent compliance. The small number of independent
scenario templates, structural recovery proxies, and pending independent
human effect-equivalence validation limit the claims; the study does not
estimate real-world prevalence or establish a general safety ranking.

% ============================================================
% GENERATIVE AI USAGE DISCLOSURE
% ============================================================

\section*{Generative AI Usage Disclosure}

Generative AI assistance was used in preparing this manuscript
for outline development, prose editing, methodological
discussion, reference verification and LaTeX/TikZ authoring.
All cited references were checked against their published
records by the authors. Any use of AI assistance in scenario
construction, implementation, annotation, analysis or figure
generation will be recorded here with the verification
procedure applied to each, in accordance with the venue's
requirements. The authors are responsible for all research
decisions, results and citations.

% ============================================================
% APPENDIX
% ============================================================

\appendix

\section{Configuration Record}
\label{app:config}


Table~\ref{tab:config} records the fields fixed before data
collection for every evaluated configuration. Reproducing a
reported rate requires all of them: tool-calling
implementations and serving backends differ in ways that change
multi-step behavior independently of the
model~\cite{chen2026tracesafe}.


\begin{table}[t]
\caption{
Fields recorded per evaluated configuration.
}
\label{tab:config}
\small
\setlength{\tabcolsep}{3pt}

\begin{tabularx}{\columnwidth}{
@{}
>{\raggedright\arraybackslash}p{2.1cm}
>{\raggedright\arraybackslash}X
@{}
}
\toprule
Field & Recorded value \\
\midrule
Model configurations & Qwen2.5-7B, Qwen2.5-14B, Mistral Small 3.2 24B \\
Mistral backend & Local Ollama 0.40.11 on A100 80 GB \\
Mistral model tag & \texttt{mistral24-detour:latest} (record blob ID) \\
Mistral quantization & Ollama distribution; verify with \texttt{ollama show} \\
Mistral context & 8192 configured tokens \\
Mistral output cap & 900 tokens \\
Audit repetitions & 3 per scenario and feedback condition \\
Clean repetitions & 2 per clean task \\
Clean-task success & 12/12 for each reported model \\
Clean chaining & 2/2 for each reported model \\
Earlier Qwen config & See their saved v8 manifests; do not assume same backend \\
\bottomrule
\end{tabularx}
\end{table}


% ============================================================
% REFERENCES
% ============================================================

\bibliographystyle{ACM-Reference-Format}
\bibliography{reference}

\end{document}